\documentclass[tikz, border=10pt]{standalone}
\usepackage{tikz-3dplot}
\usepackage{amsmath, amssymb}

\begin{document}
	
	% 设置 3D 观察视角
	\tdplotsetmaincoords{65}{110}
	
	\begin{tikzpicture}[scale=2.2, tdplot_main_coords, >=stealth, line cap=round, line join=round]
		
		% ================= 参数定义 =================
		\def\R{3.0}           % 流形曲率半径 (球体半径)
		\def\thetaP{45}       % 点 p 的极角坐标 x^1 = \theta
		\def\phiP{60}         % 点 p 的方位角坐标 x^2 = \phi
		
		% 计算自然基底的真实长度比例 (微分几何中 \partial_\phi 的范数是 R\sin\theta)
		\pgfmathsetmacro{\lenTheta}{\R * 0.4} % 缩放显示比例
		\pgfmathsetmacro{\lenPhi}{\R * sin(\thetaP) * 0.4}
		
		% 定义点 p
		\tdplotsetcoord{p}{\R}{\thetaP}{\phiP}
		
		% ================= 1. 绘制流形 M (背景网格) =================
		
		% 绘制赤道和边界圆 (Silhouette)
		\begin{scope}[tdplot_screen_coords]
			\draw[gray!70, line width=0.6pt] (0,0) circle (\R);
		\end{scope}
		
		% 绘制流形上的坐标网格 (Coordinate lines)
		% 纬度圈 (x^1 = const, x^2 变化)
		\foreach \t in {20, 45, 70} {
			\pgfmathsetmacro{\rad}{\R*sin(\t)}
			\pgfmathsetmacro{\zcoord}{\R*cos(\t)}
			% 修正点：将计算结果赋值给变量，去除多余大括号，并增加结尾分号
			\tdplotdrawarc[gray!40, thin, dashed]{(0,0,\zcoord)}{\rad}{0}{360}{}{};
		}
		
		% 经度圈 (x^2 = const, x^1 变化)
		\foreach \p in {0, 30, 60, 90, 120, 150} {
			\draw[gray!40, thin, dashed, domain=0:180, samples=40] 
			plot ({\R*sin(\x)*cos(\p)}, {\R*sin(\x)*sin(\p)}, {\R*cos(\x)});
		}
		
		% 高亮穿过 p 点的坐标曲线 (Coordinate curves generating the basis)
		% x^1 曲线 (经线, 红色虚线)
		\draw[red!60, thick, dashed, domain=0:180, samples=40] 
		plot ({\R*sin(\x)*cos(\phiP)}, {\R*sin(\x)*sin(\phiP)}, {\R*cos(\x)});
		
		% x^2 曲线 (纬线, 蓝色虚线)
		% 修正点：提取坐标 Z 轴高度和半径的数学计算，防止内部大括号解析错误
		\pgfmathsetmacro{\pz}{\R*cos(\thetaP)}
		\pgfmathsetmacro{\pr}{\R*sin(\thetaP)}
		\tdplotdrawarc[blue!60, thick, dashed]{(0,0,\pz)}{\pr}{0}{360}{}{};
		
		% ================= 2. 绘制任意光滑曲线 \gamma(\lambda) =================
		% 构造一条流形上的穿过 p 的曲线
		\draw[green!60!black, thick, domain=-40:40, samples=30] 
		plot ({\R*sin(\thetaP + 0.5*\x)*cos(\phiP + \x)}, 
		{\R*sin(\thetaP + 0.5*\x)*sin(\phiP + \x)}, 
		{\R*cos(\thetaP + 0.5*\x)});
		\node[text=green!60!black, scale=0.9, anchor=north east] at 
		({\R*sin(\thetaP-15)*cos(\phiP-30)}, {\R*sin(\thetaP-15)*sin(\phiP-30)}, {\R*cos(\thetaP-15)}) 
		{$\gamma(\lambda)$};
		
		% ================= 3. 构造局部切空间 T_p M =================
		
		% 标架变换：将局部原点平移至 p，坐标轴与自然基底对齐
		\tdplotsetrotatedcoords{\phiP}{\thetaP}{0}
		\tdplotsetrotatedcoordsorigin{(p)}
		
		\begin{scope}[tdplot_rotated_coords]
			% --- A. 绘制切空间平面 T_p M ---
			\def\planeSize{1.8}
			\filldraw[fill=yellow!20, opacity=0.4, draw=yellow!60!black, line width=0.8pt] 
			(-\planeSize, \planeSize, 0) -- (\planeSize, \planeSize, 0) -- 
			(\planeSize, -\planeSize, 0) -- (-\planeSize, -\planeSize, 0) -- cycle;
			
			\node[anchor=south west, text=yellow!40!black, scale=1.1, inner sep=4pt] 
			at (\planeSize, \planeSize, 0) {$T_p\mathcal{M}$};
			
			% --- B. 绘制自然基底 (Coordinate Basis Vectors) ---
			% 注意：在微分几何中，\partial_\theta 沿子午线向下，\partial_\phi 沿纬线方向
			
			% \partial_\theta = \partial / \partial x^1 (对应新 x 轴)
			\draw[line width=1.4pt, ->, red!80!black] (0,0,0) -- (\lenTheta, 0, 0) 
			node[below left] {$\boldsymbol{\partial}_\theta$};
			
			% \partial_\phi = \partial / \partial x^2 (对应新 y 轴)
			\draw[line width=1.4pt, ->, blue!80!black] (0,0,0) -- (0, \lenPhi, 0) 
			node[below right] {$\boldsymbol{\partial}_\phi$};
			
			% --- C. 绘制沿曲线 \gamma 的切矢量 V \in T_p M ---
			% 根据构造，\gamma 曲线的切矢量在基底上的展开 V = 0.5 \partial_\theta + 1.0 \partial_\phi
			\draw[line width=1.6pt, ->, green!50!black] (0,0,0) -- (0.5*\lenTheta, \lenPhi, 0) 
			node[above right, xshift=-2pt] {$V$};
		\end{scope}
		
		% ================= 4. 绘制点 p 与流形标记 =================
		\fill[black] (p) circle (1.5pt) node[above right, xshift=1pt, yshift=2pt] {$p$};
		\node[anchor=north west, scale=1.3] at (-2, 2.5, 0) {$\mathcal{M}$};
		
	\end{tikzpicture}
	
\end{document}